\begin{theorem}
$\sk{k} = \skp{k}$ и $\pk{k} = \pkp{k}$ (Определение классов через сертификаты и альтернирующую МТ эквивалентны).
\end{theorem}

\begin{proof} 
\begin{itemize}
    \item $\skp{k} \subset \sk{k}$. Пусть $A$ распознается альтернирующей МТ $T$, в которой число перемен состояний не более $k - 1$. Тогда она лежит в $\sk{k}$: соответствующими литералами $y_1, \ldots, y_k$ в формуле будет последовательность выборов конфигурации в каждой ветви. Пусть МТ $M$ на входе $(x, y_1, \ldots, y_k)$ моделирует работу $T$ на входе $x$ вдоль ветви $(y_1, \ldots, y_k)$. Тогда по определению альтернирующей МТ слово $x$ будет принято.
    \item $\sk{k} \subset \skp{k}$. $x \in A \Longleftrightarrow \exists y_1 \forall y_2 \exists y_3 \ldots M(x, y_1, \ldots, y_k) = 1$. Тогда альтернирующая машина сначала в $\Sigma-состоянии$ ($\exists$-состоянии) пишет биты $y_1$, затем перейдет в $\Pi$-состояние ($\forall$-состояние) и напишет биты $y_2$ и так далее. Затем запускаем детерминированное вычисление $M$.
\end{itemize}

Аналогично про классы $\pk{k}$ и $\pkp{k}$.
\end{proof}

Тогда можно определять \ph как через объединение \sk{k} (сертификатное через набор значений литералов), так и через объединение \skp{k} (альтернирующие МТ).
